% =====================================================================
%  11 -- THE ML TARGET, DECIDED  (source node S15, Part V).
%
%  Source: tier3.md lines 4300-4609, in full and in order --
%          Part V header (4300-4303), V.1 (4305-4346), V.2 (4347-4374),
%          V.3 (4375-4420), V.4 (4421-4462), V.5 (4463-4490),
%          V.5b (4491-4558), V.5c (4559-4594), V.6 self-check (4595-4609,
%          six questions -> sec:selfcheck-target).
%
%  Re-homings: none.  Cross-references to other parts of tier3.md are
%  rendered as live \S\ref{} per the label registry; no wording of a
%  cross-reference sentence had to change (all targets keep their
%  relative direction: Part IV precedes, Part VI follows).
%  Tier 2 / Tier 1 / Tier 0 pointers stay as plain text per the brief.
% =====================================================================

\section{The ML target, decided}
\label{part:target}

\textbf{A short node, because Tier~2 \S I did the algebra.} What remains is
what the verdict changed.

% =====================================================================
\subsection{The two targets, restated}
\label{sec:twotargets}

\textbf{Target A.} The network predicts a signed net per-reaction flux
(internally, in an asinh latent; the decoded quantity is linear), and the
composition change is
%
\begin{equation}
\label{eq:targeta}
\dd\mathbf Y \;=\; \nu\,\boldsymbol\varphi ,
\end{equation}
%
with $\nu$ the fixed integer stoichiometric matrix. Because $C\nu = 0$ by
construction (Tier~2 \S I.7), \textbf{every conservation law holds exactly for
any $\varphi$}, including the output of an untrained network with random
weights. Output dimension: 607 / 1518.

\textbf{Target B.} The network predicts $\dd Y$ directly in a \textbf{linear}
output space, and the result is projected onto the constraint manifold:
%
\begin{equation}
\label{eq:targetb}
P \;=\; I - C\Tr(CC\Tr)^{-1}C
\qquad\text{(shipped as the QR form } I - QQ\Tr\text{)},
\qquad \dd\mathbf Y_{\mathrm{cons}} = P\,\dd\mathbf Y_{\mathrm{raw}} .
\end{equation}
%
Output dimension: $n + 3 = 83$ / 154 (the extended vector includes $e^-$,
$\nu$, $\bar\nu$).

\textbf{The theorem that reframes the choice} (Tier~2 \S I.10): the two targets
reach \textbf{the same set}. $\image(\nut) = \operatorname{null}(C)$ exactly. So
the decision is \textbf{not about expressivity} --- it never was --- and any
document framing it that way is wrong (\foreignreg{2}{R-22} is a referee pass
to check that none does).

What the decision \emph{is} about, in three words: \textbf{parameterisation,
conditioning, supervision.}

\begingroup\small
\begin{tabularx}{\textwidth}{@{}L{3.0cm} Y Y@{}}
\toprule
 & \textbf{Target A} & \textbf{Target B} \\
\midrule
conservation & exact, structural & exact, by projection \\
output dim & 607 / 1518 & 83 / 154 \\
error attribution & per named reaction & per species \\
equilibrium structure & expressible (mask columns) & not expressible \\
conditioning & inherits $1/\kappa$ on cancelled columns & cancellation hidden
  inside $\dd Y$ labels \\
supervision available & $\Phi$ labels must be \textbf{generated}
  (\S\ref{sec:phisupervision}) & $\Delta Y$ labels ship with the dataset \\
identifiability & \textbf{528 / 1368 unidentifiable directions} under a
  $\Delta Y$-only loss (\foreignreg{2}{R-03}) & fully identified \\
\bottomrule
\end{tabularx}
\endgroup

% =====================================================================
\subsection{What the verdict decided, and what it left open}
\label{sec:decided}

\textbf{Decided: Target A, full-width.} No switch condition fires
(\S\ref{sec:whatwouldflip}). The per-reaction structure is worth keeping
precisely because the physics is concentrated in few columns --- top-20 bridges
carry 0.88 of inter-group flux (\msn{80}), top-20 weak channels carry 1.00 of
the \Ye{} budget --- and that concentration is expressible in flux space and
invisible in $\dd Y$ space.

\textbf{Decided: the equilibrium mask is gone.} The specified Component B was a
\emph{hybrid}: a Guidry physics prior, refined by a learned $L_0$ /
hard-concrete gate \citep{louizos2018}. The prior is measured empty
(\S\ref{sec:gate5}). What remains is one of:

\begin{itemize}
\item \textbf{(a) no mask} --- a full-width flux head, which is the verdict
  document's reading;
\item \textbf{(b) a purely learned gate} with no physics initialisation --- now
  an unmotivated sparsity regulariser rather than a physics-informed structure;
\item \textbf{(c) a mask against the manifold's own attractor} rather than true
  NSE --- untested, cheap, and possibly non-empty (\reg{T-04}).
\end{itemize}

The verdict picks (a) and is right to for now, but (c) has not been ruled out
and should be before Component B is written.

\textbf{Left open: the [4.0, 6.3) worst-species coverage failure}
(\S\ref{sec:gate2}), which is a \emph{supervision} problem rather than a target
problem, and whose mitigation depends on \S\ref{sec:phisupervision}.

% =====================================================================
\subsection{Why conservation must live in the decode step --- the NuGNN failure mode}
\label{sec:decodestep}

\code{CLAUDE.md} invariant \#4 says conservation lives in the decode step,
never inside latent dynamics, and that Target B's projection is valid
\textbf{only in a linear output space}. The reason is one line of algebra and
one documented failure.

A sum constraint evaluated in a warped space conserves nothing in the physical
space. If the network's output is $s = \text{signed-log}(\dd Y)$ and you project
$s$ onto $\{Cs = 0\}$, then
%
\begin{equation}
\label{eq:warpedsum}
\sum_i c_i\,s_i = 0 \quad\not\Longrightarrow\quad \sum_i c_i\,\dd Y_i = 0 ,
\end{equation}
%
because the map $s \mapsto \dd Y$ is nonlinear. The projected vector satisfies
a constraint on a transformed quantity that has no physical meaning. This is
the linear-constraint case in which architecture-level (``hard'') enforcement
is exact to machine precision \citep{beucler2021}; the linearity is the
load-bearing hypothesis.

NuGNN (\citealt{kim2026}, \arxiv{2606.04491}, \S Training Method and
\S Conclusions) documents the empirical half: with signed-log outputs, explicit
enforcement of $\sum\Delta X = 0$ --- as a loss term or an output transform,
either requiring inversion of the signed-log representation back to
$\Delta X$ --- made optimisation unstable and degraded training, so the model
shipped without an enforcing term; the paper reports that conservation was
nonetheless learned ``sufficiently well'' from data, backed by a rollout retry
(smaller $\Delta t$) when $\sum\Delta X$ ``deviated significantly from zero''.
The algebraic half --- a linear constraint imposed in warped coordinates is not
the physical constraint --- is derived here, and is not what NuGNN did (its
enforcement attempts evaluated the sum in linear $\Delta X$ space). An earlier
version of this paragraph said NuGNN ``documents exactly this instability'';
that overstated the paper (corrected in the 2026-08-18 audit). The repo encodes
the lesson as a refusal in \code{graph/projector.py}'s docstring and as
invariant \#4.

\textbf{Two corollaries that are easy to miss.}

\begin{enumerate}
\item \textbf{Target A's asinh latent is safe} precisely because the
  \emph{decoded} quantity is $\varphi$, which is linear, and the conservation
  map $\nu$ acts on $\varphi$. The nonlinearity is upstream of the conservation
  map, never between it and the output.
\item \textbf{The projector acts on the extended vector}
  $[\dd Y_{\mathrm{nuclei}}, \dd Y_{e^-}, \dd Y_{\nu}, \dd Y_{\bar\nu}]$, not
  on the nuclei alone. Projecting a nuclei-only vector against a nuclei-only
  $C$ would force $\sum Z_i\,\dd Y_i = 0$ --- i.e.\ \textbf{erase the weak
  $\dd\Ye$ signal entirely} (invariant \#3). Tier~2 \S I.9 calls this the
  laundering trap; the measured confirmation is that weak $\dd\Ye$ changes by
  $\le 3.4 \times 10^{-21}$ through projection \resultsref{2026-07-08}.
\end{enumerate}

% =====================================================================
\subsection{\texorpdfstring{$\Phi$}{Phi} supervision: proven where it matters, blocked where it does not}
\label{sec:phisupervision}

Target A predicts $\Phi$. The shipped dataset contains $\Delta Y$.
\textbf{$\Phi$ labels do not exist and must be generated}, which is what
\code{fluxes/integrate.py} (\code{ADR 0006}) is for: an augmented state
$u = [Y, \Phi]$ with $\dd\Phi_j/\dd t = R_j$, integrated by BDF
\citep{hairer1996} with a sparse analytic Jacobian, returning
$Y := Y_0 + \nu(\Phi - \Phi_0)$ so that $\Delta Y = \nu\Phi$ and conservation
hold \textbf{by construction}.

Measured \resultsref{2026-07-12}:

\begingroup\small
\begin{tabularx}{\textwidth}{@{}L{5.2cm} Y@{}}
\toprule
 & \textbf{value} \\
\midrule
internal identity residual & $2.1 \times 10^{-16}$ molar (machine precision) \\
RHS $\equiv$ engine & $\le 10^{-12}$ rel \\
BDF $\equiv$ Radau $\equiv$ BDF(rtol $10^{-10}$) & $5 \times 10^{-8}$ rel on
  the stiffest test state \\
\textbf{energy identity (invariant \#5)} & median $6 \times 10^{-6}$ \ldots\
  $2.5 \times 10^{-5}$ rel, max $8.7 \times 10^{-4}$, \textbf{fraction
  $\le 1\%$ = 1.000 in every stratum}, both nets \\
$\Delta X$ vs shipped labels, QSE window [3.3, 5.0) & median species-fraction
  in band \textbf{0.74 -- 0.99} \\
$\Delta X$ vs labels, cold strata & 0.62 -- 0.90 \\
$\Tn \ge 5$ & \textbf{unmeasurable within deadline AND label-pathological} \\
\bottomrule
\end{tabularx}
\endgroup

\textbf{Feasibility, which is the binding constraint:} median wall $\ge 383$~s
(\msn{80}) / 1519~s (\msn{151}) per state, with 17/36 and 10/24 states
\textbf{censored} at the 4800-s deadline --- and every censored state is
$\Tn \ge 5$. That is $\le 9.4$ / 2.4 states per hour per core, so the full
corpus would cost \textbf{$\ge$ 110,647 / 439,302 core-hours}. Full-corpus
$\Phi$ supervision is \textbf{infeasible}. Stratified subsets of
$10^3$--$10^4$ states in the $\Tn < 5$ band, QSE window first, cost
$\approx 4$--40 core-days and are feasible.

\textbf{The cost cause is diagnosed, which turns a wall into a task.} The
deliberately neglected tabular-EC $\rho\Ye$ Jacobian chain dominates the
weak-drift phase at $\Tn \ge 5$: the solver is taking tiny steps because it has
no derivative information for the term that is actually driving the evolution.
\code{ADR 0006}'s contingency --- add the \Ye-chain term to the Jacobian --- is
the identified unlock.

So the picture for Phase 1 is coherent:

\begingroup\small
\begin{tabularx}{\textwidth}{@{}L{1.5cm} L{3.4cm} L{2.4cm} L{3.0cm} Y@{}}
\toprule
\textbf{band} & \textbf{$\Phi$ labels} & \textbf{label physics} &
  \textbf{kill-test gates} & \textbf{verdict} \\
\midrule
\textbf{$\Tn < 5$} & generatable --- 4--40 core-days for $10^3$--$10^4$
  states & sound & all pass at median coverage & \textbf{this is where
  Target A is trained and validated} \\
\textbf{$\Tn \ge 5$} & blocked on integrator cost (unlock identified) &
  bug-displaced (\S\ref{part:pathology}) & median coverage fails in
  [5.0, 6.3) & deferred; declare a validity domain (\reg{T-03}) \\
\bottomrule
\end{tabularx}
\endgroup

% =====================================================================
\subsection{The identifiability cost nobody has paid attention to}
\label{sec:identifiability}

Tier~2 \S IV.4b derives it and it belongs in the target decision.

Under a \textbf{$\Delta Y$-only} loss, the flux head is unidentifiable in
\textbf{528 / 1368 directions} --- the dimension of $\operatorname{null}(\nu)$
restricted to the relevant column space. Any $\varphi$ and $\varphi + n$ with
$n \in \operatorname{null}(\nu)$ produce identical $\Delta Y$, identical
conservation, identical energy (if $Q \perp \operatorname{null}(\nu)$, which
Tier~2 \S IV.8b measured to be \emph{false} at the 0.55~keV rms level, but
nearly true). The loss cannot distinguish them.

Three consequences:

\begin{enumerate}
\item \textbf{The learned fluxes are not the physical fluxes}, even for a
  perfectly accurate model, unless $\Phi$-level supervision pins the null
  directions. The ``physical interpretability of errors'' advantage of
  Target A (Tier~2 \S 0.1.3) is \emph{conditional on $\Phi$ supervision}.
\item \textbf{Gradient descent will wander in the null space}, which is a
  training pathology (no gradient signal, but also no penalty) rather than an
  accuracy one.
\item \textbf{It is a second, independent cost of Target A}, alongside the
  $1/\kappa$ conditioning of \S\ref{sec:gate2} --- and it is one Target B does
  not have.
\end{enumerate}

None of that overturns the verdict; Target B's costs (loss of per-reaction
structure, loss of the equilibrium expressibility, the same label-level
cancellation) are judged larger. But the honest statement of the decision is
``Target A, with two known costs and a mitigation plan for one of them'', not
``Target A wins''.

% =====================================================================
\subsection{The output space is a simplex, and nothing in the project treats it as one}
\label{sec:simplex}

A prerequisite that no tier covers: \textbf{the geometry of the space the
answer lives in.}

A composition is a point on the $(n-1)$-simplex
%
\begin{equation}
\label{eq:simplex}
\Delta^{n-1} \;=\; \Big\{\,X \in \mathbb R^{n} \;:\; X_i \ge 0,\
\textstyle\sum_i X_i = 1 \,\Big\},
\end{equation}
%
and increments live in its tangent space, which is where Tier~2 \S I.7b's
positivity discussion starts. The project handles the \textbf{equality}
constraint beautifully --- it is the whole conservation story --- and handles
the \textbf{inequality} constraint nowhere. Three things follow that are not in
any register.

\textbf{(i) There is a mature literature for exactly this and it is uncited.}
Compositional data analysis (\citealt{aitchison1982, aitchison1986}; ilr:
\citealt{egozcue2003}) treats data on a simplex through log-ratio
transforms --- additive (alr), centred (clr) and isometric (ilr) --- under which
the simplex becomes a real vector space with its own inner product, and
standard statistics apply. The relevant facts:

\begin{itemize}
\item \textbf{clr and ilr are linear-in-log-ratio}, which is \emph{not} the
  same as ``log space'', and the distinction is exactly the one invariant \#4
  is about. A clr coordinate is $\Sigma$-free by construction: the closure
  constraint is absorbed into the coordinate system rather than imposed on top
  of it.
\item The NNN's log-softmax output (\S\ref{sec:nnnis};
  \citealt[\S 2.3--2.4]{grichener2025}) is precisely the \textbf{inverse
  clr/alr map} applied without anyone saying so --- softmax is the standard
  simplex parameterisation. So the baseline is already doing compositional data
  analysis by accident, and the vocabulary would have told it that the
  corresponding \emph{loss} should be a log-ratio distance rather than an L1 on
  logs.
\item The Aitchison geometry's pathology is also known: \textbf{zeros are not
  representable} ($\log 0 = -\infty$), and compositional data analysis has a
  whole sub-literature on zero replacement \citep[e.g.][]{martinfernandez2003}.
  This project's data has a hard floor at $10^{-15}$ (\S\ref{sec:rowisnot}) ---
  i.e.\ it has \emph{already made} a zero-replacement choice, upstream, without
  calling it that.
\end{itemize}

\textbf{(ii) Target A sidesteps the simplex and pays for it elsewhere.}
$\dd Y = \nu\varphi$ guarantees $\sum A\,\dd Y = 0$ exactly, so the
\emph{hyperplane} is respected; the \emph{facets} are not. Nothing prevents an
individual $Y_i$ from going negative, and Tier~2 \S I.7b argues it is not a
small problem here because trace species span fifteen decades. The available
options, none of which is in the plan:

\begingroup\small
\begin{tabularx}{\textwidth}{@{}L{6.4cm} Y@{}}
\toprule
\textbf{option} & \textbf{cost} \\
\midrule
diagnose only (count violations, magnitude) & free; \textbf{should exist
  already} (\foreignreg{2}{R-04}) \\
clip and renormalise & breaks conservation --- the exact thing the design
  forbids \\
project onto the simplex (Euclidean or Aitchison) & conservation-preserving if
  done in the tangent space; changes the loss landscape \\
parameterise so positivity is structural (predict in a log-ratio coordinate
  and decode) & \textbf{puts a nonlinearity between the conservation map and
  the output --- forbidden by invariant \#4} \\
\bottomrule
\end{tabularx}
\endgroup

The last row is the interesting one: \textbf{positivity and exact linear
conservation pull in opposite directions.} Structural positivity wants an
exponential map; structural conservation wants linearity. You can have either
structurally, or one structurally and the other by projection, but the obvious
way to get both at once is exactly the NuGNN failure mode. That tension is
real, is not written down anywhere in the project, and is the correct framing
of \foreignreg{2}{R-04}.

\textbf{(iii) The metric on the correction is a choice nobody recorded.}
\foreignreg{2}{R-17} notes that the Target B projector is orthogonal in the
\emph{unweighted Euclidean} metric --- the minimum-$\sum(\delta Y)^2$
correction. On a simplex, the natural metrics are Aitchison's or an
abundance-weighted one, and the difference matters precisely for trace species:
an unweighted projector distributes the correction uniformly, so it makes a
relatively enormous change to a $10^{-12}$ species and a negligible one to a
0.5 species. Registered as \reg{T-30}, extending \foreignreg{2}{R-17} with the
reason the choice is not innocuous.

% =====================================================================
\subsection{Gradients through the conservation map --- two different questions}
\label{sec:gradients}

\tierone{VIII.E.5} flags ``gradients through $\nu$ and $P$'' as a gap. It is
worth separating two things that the single flag conflates, because one is a
training question and one is the deployment requirement of
\S\ref{sec:callsite} item~2.

\textbf{Question 1 --- training.} Backpropagating through $\dd Y = \nu\varphi$
is a constant linear map, so the gradient is $\nu\Tr$, which is exact, cheap
and well conditioned ($\cond(\nu) = 41.6$ / 57.9). The real subtlety is
$\nu$'s \textbf{nullity}: the gauge freedom of Tier~2 \S I.2.1 means gradient
descent receives no signal in $\operatorname{null}(\nu)$, which is the
identifiability statement of \S\ref{sec:identifiability} seen from the
optimiser's side. Nothing diverges; the parameters simply drift in those
directions, which is a slow pathology and an unexamined one.

\textbf{Question 2 --- deployment.} The host needs $\partial\enuc/\partial\ln T$
and $\partial\enuc/\partial\ln\rho$ (\S\ref{sec:callsite} item~2, \reg{T-25}).
Through the flux route,
%
\begin{equation}
\label{eq:denucdlnT}
e_{\mathrm{nuc}} = \sum_j Q_j \Phi_j
\qquad\Longrightarrow\qquad
\frac{\partial e_{\mathrm{nuc}}}{\partial \ln T}
= \sum_j Q_j\,\frac{\partial \Phi_j}{\partial \ln T}
\;+\;\sum_j \Phi_j\,\frac{\partial Q_j}{\partial \ln T},
\end{equation}
%
and the second term is \textbf{exactly the $Q(T)$ correction that
\tierone{VIII.E.2} says is specified, unimplemented, and invisible to the
current energy gate} --- where it reaches $-19.2\%$ on
\nuc{Ca}{40}($\alpha,\gamma$)\nuc{Ti}{44} at $\Tn = 7.9$
\derivedthere[tier1 \S VIII.E.2 table]. So an open Tier-1 item and an
unrecognised deployment requirement are the same term. Closing \S VIII.E.2 is
therefore worth more than it looked: it is not only an energy-accuracy item, it
is a \emph{prerequisite for supplying the host with a correct derivative}.

That connection is new here and it is the kind of thing a cross-tier audit is
for. Folded into \reg{T-25}.

% =====================================================================
\subsection{Self-check}
\label{sec:selfcheck-target}

\begin{selfcheck}
\item % source V.6 Q1
  State the reachability theorem and say what it removes from the A-vs-B
  argument.
\item % source V.6 Q2
  Give three axes on which A and B genuinely differ, with the measured number
  supporting each.
\item % source V.6 Q3
  Why is Target A's asinh latent safe when NuGNN's signed-log conservation was
  not?
\item % source V.6 Q4
  What does projecting a \emph{nuclei-only} $\dd Y$ vector against a
  nuclei-only $C$ do to \Ye, and which invariant forbids it?
\item % source V.6 Q5
  $\Phi$ labels cost $\ge 110$k core-hours for the full corpus. What is the
  diagnosed cause, what is the unlock, and what is the feasible fallback?
\item % source V.6 Q6
  How many directions is the flux head unidentifiable in, under what loss, and
  what fixes it?
\end{selfcheck}
